\section{Literature review}
\label{sec:literature review}


% The problem of chess puzzle generation has been studied suprisingly quite extensively despite it's nieche use cases \textcolor{red}{add a bunch of references}. Chess platforms like Chess.com and Lichess generate puzzles to their users by filtering a huge amount of positions from games played on their websites. More effort in the scientific commmunity has been put into directly generating chess positions with desired properties.

% In this study, we largely follow the work of \cite{feng2025generatingcreativechesspuzzles} and train a diffusion model to generate chess positions. We use chess puzzles as a well-defined playground to study the properties of diffusion models in the field of language modelling. In particular, we attempt to improve the creative qualitites of a language model that are hard to measure with traditional metrics. In addition, we investigate methods of post-training diffusion models with reinforcement learning to improve the quality of the generated positions.

In this study, we largely follow the work of \citeauthor{feng2025generatingcreativechesspuzzles} in \cite{feng2025generatingcreativechesspuzzles} and train a diffusion model to generate chess positions. In \cite{feng2025generatingcreativechesspuzzles}, the authors first train four separate deep learning models with different architectures using supervised learning. After pretraining the models, the authors select an autoregressive model with the same architecture as modern LLMs, which performs well after the pretraining phase. They fine-tune the autoregressive model using reinforcement learning to improve the quality of the generated positions with the goal of generating positions that a human would consider creative. \citeauthor{feng2025generatingcreativechesspuzzles} find that reinforcement learning improves the quality and quantity of puzzles a lot. In the end, their model is able to generate positions that are almost on par with award winning puzzles when compared by master-level chess players and puzzle composers.

\textcolor{red}{Perhaps talk about more traditional (evolution-based) algorithms for generating chess puzzles. (measure of machine intelligence)}


We use chess puzzle generation as a well-defined playground to test and improve the capabilities of reinforcement learning with masked diffusion language models. Originally developed by \citeauthor{austin2021structureddenoisingdiffusionmodels} in \cite{austin2021structureddenoisingdiffusionmodels} and later scaled up by many others, see e.g. \cite{nie2025largelanguagediffusionmodels, ye2025dream7bdiffusionlarge}, masked diffusion models have shown promise in the language modelling domain. In particular, they offer benefits in terms of efficiency during inference, as one can vary the amount of inference time compute depending on the task by using different amount of denoising steps. However, reinforcement learning, in particular reinforcement learning from human feedback (RLHF) and reinforcement learning with verifiable rewards (RLVR), is a huge part of modern LLM training \cite{shao2024deepseekmathpushinglimitsmathematical, ouyang2022traininglanguagemodelstofollowinstructionswithhumanfeedback, lambert2025tulu}, but is not as well-studied in the context of discrete diffusion models. \citeauthor{black2024trainingdiffusionmodelsreinforcement} provide a framework for reinforcement learning with diffusion models in general, but for language modeling, the research is more limited. Due to computational constraints, the same approach is not possible for discrete diffusion language models, however, methods have been proposed to approximate the likelihood of individual samples \cite{ou2025principledrldiffusionllms, zhao2025d1scalingreasoningdiffusion,tang2026wd1weightedpolicyoptimization,rojas2025improvingreasoningdiffusionlanguage}. We use and RL algorithm inspired by \cite{ou2025principledrldiffusionllms} and \cite{rojas2025improvingreasoningdiffusionlanguage} to our chess puzzle generation task.
